\section{Discussion}\label{s:discussion}

The main finding is that resistance affected information gathering before it affected the final diagnosis. Question responsiveness and information yield decreased steadily, while diagnostic classification remained largely accurate from low through high intensity. Classification accuracy fell sharply only under extreme resistance. The category results also showed different forms of disruption, including short or absent answers, irrelevant detail, challenges to the assessment, and attempts to control its direction.

\subsection{Conversation Quality and Client-Side Measures}

Resistance affected the patient-response measures more strongly than the broad conversation rubric. Responsiveness and information yield fell steadily from baseline through extreme, while the rubric mean decreased more modestly. The rubric combines several properties of the conversation, whereas responsiveness and information yield focus on whether each patient answer addresses the current question and provides useful evidence. A conversation can therefore complete the questionnaire and retain a moderate rubric score even when the patient supplies much less diagnostic information. In a questionnaire-based assessment, completing the interview and gathering sufficient information are not the same outcome.

Silence and vague answers reduced the amount of text, while irrelevant verbosity greatly increased it. Contesting and controlling often produced answers close to the baseline length but avoided the requested information or interfered with the assessment. Responsiveness and information yield captured these less visible forms of disruption.

\subsection{Diagnostic Performance under Reduced Information}

Diagnostic accuracy changed little from low through high intensity, despite lower responsiveness and information yield. These conversations still contained non-activated turns that could disclose profile information. In addition, the clinical profile remained in the patient agent's prompt on activated turns, so resistant responses could still reveal parts of the case. Either route may have left enough evidence in the dialogue for the diagnostic agent to identify the assigned disorder. The present analysis does not trace which parts of the dialogue supported each diagnosis, so it cannot determine which route mattered most.

Diagnostic certainty changed earlier than accuracy. From low to high intensity, certainty decreased even though classification accuracy remained high. Resistance therefore affected how confidently the reports stated their conclusions before it affected whether the primary diagnosis was correct.

Under extreme resistance, classification accuracy fell from 0.963 at high intensity to 0.583. The relationship between information yield and accuracy was statistically reliable when all conditions were analyzed together, but not after extreme resistance was removed. Among the 925 baseline-to-high conversations, 908 were classified correctly, leaving only 17 errors for comparison. This ceiling effect makes smaller changes difficult to detect. The experiment therefore shows a clear failure under extreme resistance, but it does not locate the amount of information at which diagnostic accuracy begins to decline.

Reframing requires separate interpretation. Its truth map controls whether the reframing instruction can be added to a turn, but it does not remove the clinical profile from the patient prompt. On ineligible turns, the patient returns to an ordinary profile-grounded response and may still disclose symptoms or background information from that profile. This additional information, together with the lower activation rate, may explain why reframing retained high diagnostic accuracy. Its result does not show that reframing is inherently less disruptive. Future comparisons should report activation among eligible turns and measure the information disclosed on non-activated turns.

After reframing was excluded, classification accuracy did not differ across the remaining categories, but diagnostic justification and certainty did. Resistance category therefore affected how the report supported and expressed its conclusion more clearly than whether it selected the correct primary label. This is another reason not to evaluate diagnostic performance through accuracy alone.

The extreme condition also exposed a problem in how evidence was attributed. Some reports gave a specific diagnosis despite receiving little usable information. In one inspected example, the report treated the wording of an unanswered questionnaire item as if the patient had endorsed the symptom. Because the questionnaire itself names possible symptoms, the diagnostic agent must distinguish what was asked from what the patient confirmed. The reviewed examples are not a systematic error analysis, but they motivate two safeguards: linking each claimed symptom to a patient response and allowing an ``insufficient information'' outcome when the dialogue does not support a diagnosis.

\subsection{Limitations and Future Work}

The dataset consists entirely of synthetic conversations generated from five fixed clinical profiles. The repeated runs reuse the same profile and condition prompt, so they capture variation in model generation and resistance activation rather than variation between patient cases. Each disorder is also represented by only one profile. The results therefore cannot show whether some disorders are inherently more sensitive to resistance than others. A broader study would require multiple profiles for each disorder and clinician review of the simulated cases.

The intensity conditions are not four equally spaced points on one scale. Low, medium, and high use probabilistic activation, while extreme activates every eligible turn and always uses Level 3 guidance. Extreme should therefore be read as a stress test rather than as the next step above high intensity. Additional probabilistic conditions between high and extreme would help locate where diagnostic accuracy begins to decline. Reframing also differs from the other categories because its truth map limits resistance to profile-positive items. On other turns, the patient may still disclose information from the clinical profile. Moreover, the truth maps were constructed for the five profiles in this study and were not independently validated by a clinical expert. Reframing should therefore be evaluated separately, with activation reported over eligible turns and information disclosure measured on the remaining turns.

Gemini 2.5 Flash is used for all three agents and for most of the automatic evaluation. Other models may follow the patient prompts differently, so the same resistance profiles may not produce identical results. An evaluator based on the same model may also favor outputs produced by that model. Resistance fidelity was not checked by multiple human annotators, and the inspected diagnostic errors were selected as examples rather than systematically coded. These results describe the behavior of the current simulation, not its clinical validity. Replication with other models, human ratings, and systematic coding of diagnostic errors is needed.

%%% Local Variables:
%%% mode: latex
%%% TeX-master: "../thesis"
%%% End:
